\documentclass[12pt]{article}
\usepackage[T1]{fontenc}
\usepackage[utf8]{inputenc}
\usepackage{lmodern}
\usepackage{textcomp}
\usepackage{enumitem}
\usepackage{geometry}
\usepackage{graphicx}
\usepackage{amsmath, amssymb}
\geometry{margin=1in}
\begin{document}
\begin{center}
Physics - Undergraduate Level Assessment 17 \
Total Marks: 40 \
Answer all questions.
\end{center}

\section*{Multiple Choice (10 marks)}
\begin{enumerate}[label=\arabic*.]
\item If the absolute temperature of a blackbody is increased by a factor of 3, the energy radiated per second per unit area does which of the following?
\begin{enumerate}[label=(\alph*)]
    \item Decreases by a factor of 81.
    \item Decreases by a factor of 9.
    \item Increases by a factor of 9.
    \item Increases by a factor of 81.
\end{enumerate}
\item Which of the following is true about any system that undergoes a reversible thermodynamic process?
\begin{enumerate}[label=(\alph*)]
    \item There are no changes in the internal energy of the system.
    \item The temperature of the system remains constant during the process.
    \item The entropy of the system and its environment remains unchanged.
    \item The entropy of the system and its environment must increase.
\end{enumerate}
\item An atom has filled n = 1 and n = 2 levels. How many electrons does the atom have?
\begin{enumerate}[label=(\alph*)]
    \item 2
    \item 4
    \item 6
    \item 10
\end{enumerate}
\item If the total energy of a particle of mass m is equal to twice its rest energy, then the magnitude of the particle's relativistic momentum is
\begin{enumerate}[label=(\alph*)]
    \item mc/2
    \item mc/($2^(1$/2))
    \item mc
    \item ($3^(1$/2))mc
\end{enumerate}
\item The rest mass of a particle with total energy 5.0 GeV and momentum 4.9 GeV/c is approximately
\begin{enumerate}[label=(\alph*)]
    \item 0.1 GeV/$c^2$
    \item 0.2 GeV/$c^2$
    \item 0.5 GeV/$c^2$
    \item 1.0 GeV/$c^2$
\end{enumerate}
\end{enumerate}

\section*{True / False (10 marks)}
\textit{Answer True or False}
\begin{enumerate}[label=\arabic*.]
\setcounter{enumi}{5}
\item True or False: The attraction between Cooper pairs in superconductors is primarily due to the strong nuclear force.
\item True or False: Electromagnetic radiation emitted from a nucleus is most commonly in the form of gamma rays.
\item True or False: In a single-electron atom with l = 2, there are only 3 allowed values for the magnetic quantum number m\_l.
\item True or False: The negative muon, mu^-, has properties most similar to mesons, which are composite particles made of quarks.
\item True or False: Gas lasers utilize transitions that involve the energy levels of free atoms to produce laser light.
\end{enumerate}

\section*{Long Form Response (20 marks)}
\textit{Answer each question with a clear and complete explanation.}
\begin{enumerate}[label=\arabic*.]
\setcounter{enumi}{10}
\item Discuss how the internal energy of an ideal gas changes during a constant volume process, specifically analyzing the relationship between heat added and internal energy increase.
\item Discuss the allowed transitions for an electron in the n = 4, l = 1 state of hydrogen, and analyze which final states are not achievable, providing reasoning for your conclusions.
\end{enumerate}

\vfill
\noindent\textit{End of Paper}
\end{document}